% Versión LaTeX de docs/p1-datos.md (2026-11-14). La fuente es el .md: si cambia, actualizar aquí.
\documentclass[11pt,a4paper]{article}
% Preámbulo común de la documentación LaTeX de TerraPulse AI.
% Se incluye con % Preámbulo común de la documentación LaTeX de TerraPulse AI.
% Se incluye con % Preámbulo común de la documentación LaTeX de TerraPulse AI.
% Se incluye con \input{preambulo} desde cada documento. Compilar con pdflatex (o latexmk -pdf).
\usepackage[utf8]{inputenc}
\usepackage[T1]{fontenc}
\usepackage{lmodern}
\usepackage[spanish,es-noshorthands,es-tabla,es-nodecimaldot]{babel}
\usepackage{amsmath,amssymb}
\usepackage[margin=2.3cm]{geometry}
\usepackage{array,booktabs,longtable}
\usepackage[normalem]{ulem}
\usepackage{xcolor}
\usepackage{enumitem}
\usepackage{listings}
\usepackage[hidelinks]{hyperref}

\setlength{\parindent}{0pt}
\setlength{\parskip}{0.5em}
\setlist{itemsep=2pt,topsep=2pt}
\renewcommand{\arraystretch}{1.25}
\Urlmuskip=0mu plus 1mu

% Símbolos frecuentes
\newcommand{\sz}{\ensuremath{\sigma^0}}
\newcommand{\gz}{\ensuremath{\gamma^0}}
\newcommand{\bz}{\ensuremath{\beta^0}}
\newcommand{\dsz}{\ensuremath{\Delta\sigma^0}}
\newcommand{\ok}{\ensuremath{\checkmark}}
\newcommand{\flecha}{\ensuremath{\rightarrow}}
\newcommand{\aprox}{\ensuremath{\approx}}
\newcommand{\por}{\ensuremath{\times}}

% Recuadro de nota (equivale a los ">" del Markdown)
\newenvironment{nota}
  {\begin{quote}\small\itshape}
  {\end{quote}}

% Bloques de código y comandos
\lstset{
  basicstyle=\ttfamily\small,
  breaklines=true,
  columns=fullflexible,
  keepspaces=true,
  frame=single,
  rulecolor=\color{black!30},
  backgroundcolor=\color{black!3},
  xleftmargin=4pt, xrightmargin=4pt,
  extendedchars=true,
  literate={á}{{\'a}}1 {é}{{\'e}}1 {í}{{\'i}}1 {ó}{{\'o}}1 {ú}{{\'u}}1
           {Á}{{\'A}}1 {É}{{\'E}}1 {Í}{{\'I}}1 {Ó}{{\'O}}1 {Ú}{{\'U}}1
           {ñ}{{\~n}}1 {Ñ}{{\~N}}1
}
 desde cada documento. Compilar con pdflatex (o latexmk -pdf).
\usepackage[utf8]{inputenc}
\usepackage[T1]{fontenc}
\usepackage{lmodern}
\usepackage[spanish,es-noshorthands,es-tabla,es-nodecimaldot]{babel}
\usepackage{amsmath,amssymb}
\usepackage[margin=2.3cm]{geometry}
\usepackage{array,booktabs,longtable}
\usepackage[normalem]{ulem}
\usepackage{xcolor}
\usepackage{enumitem}
\usepackage{listings}
\usepackage[hidelinks]{hyperref}

\setlength{\parindent}{0pt}
\setlength{\parskip}{0.5em}
\setlist{itemsep=2pt,topsep=2pt}
\renewcommand{\arraystretch}{1.25}
\Urlmuskip=0mu plus 1mu

% Símbolos frecuentes
\newcommand{\sz}{\ensuremath{\sigma^0}}
\newcommand{\gz}{\ensuremath{\gamma^0}}
\newcommand{\bz}{\ensuremath{\beta^0}}
\newcommand{\dsz}{\ensuremath{\Delta\sigma^0}}
\newcommand{\ok}{\ensuremath{\checkmark}}
\newcommand{\flecha}{\ensuremath{\rightarrow}}
\newcommand{\aprox}{\ensuremath{\approx}}
\newcommand{\por}{\ensuremath{\times}}

% Recuadro de nota (equivale a los ">" del Markdown)
\newenvironment{nota}
  {\begin{quote}\small\itshape}
  {\end{quote}}

% Bloques de código y comandos
\lstset{
  basicstyle=\ttfamily\small,
  breaklines=true,
  columns=fullflexible,
  keepspaces=true,
  frame=single,
  rulecolor=\color{black!30},
  backgroundcolor=\color{black!3},
  xleftmargin=4pt, xrightmargin=4pt,
  extendedchars=true,
  literate={á}{{\'a}}1 {é}{{\'e}}1 {í}{{\'i}}1 {ó}{{\'o}}1 {ú}{{\'u}}1
           {Á}{{\'A}}1 {É}{{\'E}}1 {Í}{{\'I}}1 {Ó}{{\'O}}1 {Ú}{{\'U}}1
           {ñ}{{\~n}}1 {Ñ}{{\~N}}1
}
 desde cada documento. Compilar con pdflatex (o latexmk -pdf).
\usepackage[utf8]{inputenc}
\usepackage[T1]{fontenc}
\usepackage{lmodern}
\usepackage[spanish,es-noshorthands,es-tabla,es-nodecimaldot]{babel}
\usepackage{amsmath,amssymb}
\usepackage[margin=2.3cm]{geometry}
\usepackage{array,booktabs,longtable}
\usepackage[normalem]{ulem}
\usepackage{xcolor}
\usepackage{enumitem}
\usepackage{listings}
\usepackage[hidelinks]{hyperref}

\setlength{\parindent}{0pt}
\setlength{\parskip}{0.5em}
\setlist{itemsep=2pt,topsep=2pt}
\renewcommand{\arraystretch}{1.25}
\Urlmuskip=0mu plus 1mu

% Símbolos frecuentes
\newcommand{\sz}{\ensuremath{\sigma^0}}
\newcommand{\gz}{\ensuremath{\gamma^0}}
\newcommand{\bz}{\ensuremath{\beta^0}}
\newcommand{\dsz}{\ensuremath{\Delta\sigma^0}}
\newcommand{\ok}{\ensuremath{\checkmark}}
\newcommand{\flecha}{\ensuremath{\rightarrow}}
\newcommand{\aprox}{\ensuremath{\approx}}
\newcommand{\por}{\ensuremath{\times}}

% Recuadro de nota (equivale a los ">" del Markdown)
\newenvironment{nota}
  {\begin{quote}\small\itshape}
  {\end{quote}}

% Bloques de código y comandos
\lstset{
  basicstyle=\ttfamily\small,
  breaklines=true,
  columns=fullflexible,
  keepspaces=true,
  frame=single,
  rulecolor=\color{black!30},
  backgroundcolor=\color{black!3},
  xleftmargin=4pt, xrightmargin=4pt,
  extendedchars=true,
  literate={á}{{\'a}}1 {é}{{\'e}}1 {í}{{\'i}}1 {ó}{{\'o}}1 {ú}{{\'u}}1
           {Á}{{\'A}}1 {É}{{\'E}}1 {Í}{{\'I}}1 {Ó}{{\'O}}1 {Ú}{{\'U}}1
           {ñ}{{\~n}}1 {Ñ}{{\~N}}1
}


\title{P1 --- Datos y adquisición}
\author{TerraPulse AI}
\date{Cierre de la Fase 1, 14 de noviembre de 2026}

\begin{document}
\maketitle

Documentación del trabajo de P1 al cierre de la Fase 1 (2026-11-14). Reglas del proyecto en
\texttt{AGENTS.md}; plan y roles en \emph{Plan de trabajo} (\texttt{plan-de-trabajo.tex}); resultados
científicos y fuentes de validación en \emph{Metodología} (\texttt{metodologia.tex}).

\section{Responsabilidad}

Conseguir recortes NISAR reales de la ROI con sus metadatos, y garantizar que:

\begin{itemize}
  \item cada recorte se puede rastrear hasta la escena, la URL y la versión del código que lo produjeron;
  \item el archivo descargado es la escena que dice el catálogo;
  \item las dos fechas de un par son comparables y comparten la misma malla;
  \item se conoce qué parte de la ROI cubre cada escena.
\end{itemize}

P1 \textbf{no} convierte a dB, no aplica máscaras ni calcula cambios: eso es de P2. Los recortes son un espejo fiel del producto.

\section{Estado}

\begin{longtable}{p{0.48\linewidth}p{0.45\linewidth}}
\toprule
\textbf{Criterio de la Fase 1} & \textbf{Estado} \\
\midrule
\endhead
\path{fetch_subset} funcionando con \texttt{.netrc}, solo Frequency A & \ok{} Probado en remoto y en local \\
Dos recortes que cumplen el contrato C1 en \path{data/raw/} & \ok{} Par T62/F10 \\
Misma malla confirmada & \ok{} \\
Cobertura de la ROI medida & \ok{} 99.57\,\% y 99.61\,\% \\
Par anotado en \path{config.yaml} & \ok{} \\
Pruebas & \ok{} 44 casos pasan con \texttt{pytest} (ver §\ref{sec:pruebas}) \\
Mínimo de cobertura aceptable & \textit{Pendiente:} lo fija P3 \\
\bottomrule
\end{longtable}

\section{Archivos de P1}

\begin{longtable}{p{0.36\linewidth}p{0.57\linewidth}}
\toprule
\textbf{Archivo} & \textbf{Qué hace} \\
\midrule
\endhead
\path{pipeline/config.yaml} & ROI, fuente, producto, criterios de comparabilidad y par de trabajo (\texttt{pair}) \\
\path{pipeline/sources/base.py} & Contrato de una fuente: \texttt{DataSource} (\texttt{search}, \path{fetch_subset}, \path{fetch_subsets}) y \texttt{SceneMeta} \\
\path{pipeline/sources/asf.py} & Fuente ASF: búsqueda con \path{asf_search}, inicio de sesión en Earthdata y apertura remota o local del \texttt{.h5} \\
\path{pipeline/sources/subset.py} & Recorte independiente de la fuente: identidad, ventana, GeoTIFF, cobertura y \path{scene.json} \\
\path{pipeline/sources/__init__.py} & \path{get_source(name)}: elige la fuente según \texttt{source:} en \path{config.yaml} \\
\path{pipeline/catalog.py} & CLI: consulta la fuente, agrupa en pilas y guarda el catálogo en JSON \\
\path{pipeline/download.py} & CLI: recorta escenas del catálogo y compara el par \\
\path{pipeline/nisar_h5.py} & Lectura por ventanas de los \texttt{.h5} NISAR (malla, ventana de la ROI, lectura por franjas, dB, multilook) \\
\path{pipeline/etapa1.py}, \path{etapa1_figuras.py} & CLI de inspección de una escena local (figuras y valores de control) \\
\path{pipeline/tests/test_sources.py}, \path{test_nisar_h5.py}, \path{test_subset.py} & Pruebas con datos sintéticos \\
\bottomrule
\end{longtable}

\path{pipeline/compat.py} (comparabilidad) lo usa P1, pero es de P3.

\section{Requisitos}

\begin{enumerate}
  \item Entorno \texttt{nisar} (ver README §3).
  \item Cuenta NASA Earthdata con sus credenciales en \textbf{una} de estas formas (README §7):
  \begin{itemize}
    \item \path{EARTHDATA_USERNAME} y \path{EARTHDATA_PASSWORD} en el \texttt{.env} de la raíz del repo o como variables de
          entorno. \textbf{Tienen prioridad} y son la opción para el servidor; las variables del sistema no se sobrescriben con el archivo.
    \item \texttt{\%USERPROFILE\%\textbackslash .netrc} con
          \texttt{machine urs.earthdata.nasa.gov login TU\_USUARIO password TU\_CONTRASEÑA}.
          En Windows, \texttt{earthaccess} busca \path{_netrc}; el pipeline usa \texttt{.netrc} si \path{_netrc} no existe.
  \end{itemize}
  \item Haber iniciado sesión una vez en ASF Vertex y aceptado sus términos.
\end{enumerate}

Los comandos se ejecutan \textbf{desde la raíz del repo} y con \texttt{python -m}:

\begin{lstlisting}
conda activate nisar
cd C:\terrapulse
\end{lstlisting}

\subsection*{Obtener los datos después de clonar}

\path{data/} no está en git, así que un repo recién clonado no tiene recortes. Los datos se generan desde el servidor de ASF:

\begin{lstlisting}
conda activate nisar
cd C:\terrapulse
python -m pipeline.catalog
python -m pipeline.download NISAR_L2_PR_GCOV_023_062_A_010_4005_DHDH_A_20260618T114010_20260618T114044_P05023_N_F_J_001 NISAR_L2_PR_GCOV_033_062_A_010_4005_DHDH_A_20261016T114006_20261016T114041_P05023_N_F_J_001
\end{lstlisting}

\begin{itemize}
  \item \path{pipeline.catalog} va \textbf{primero}: \texttt{download} busca las escenas en \path{data/catalog/}. Sin catálogo falla con «No hay catálogos».
  \item \path{pipeline.download} lee el \texttt{.h5} \textbf{directamente del servidor} y baja solo los bloques de la ROI
        (\textasciitilde 230\,MB por escena). \textbf{No hace falta descargar el \texttt{.h5} completo} (\textasciitilde 7\,GB).
  \item Resultado: \path{data/raw/<scene_id>/} con 5 \texttt{.tif} y \path{scene.json} por escena. Tarda \textasciitilde 3--4 minutos y debe
        terminar con \texttt{COMPARABLE}. Los \path{scene_id} del par también están en \path{config.yaml -> pair}.
  \item Los valores de los píxeles, la malla y la cobertura salen idénticos para todo el equipo. En \path{scene.json} solo cambian la
        fecha de creación y, si se trabaja en otro commit, el \path{code_version}.
  \item Para que \path{code_version} no termine en \texttt{-dirty}, ejecutarlo con el árbol de git limpio (\texttt{git status -{}-short} sin salida).
  \item Si ASF reprocesa una escena, su \path{scene_id} cambia y el comando no la encontrará: revisar el catálogo nuevo y no mezclar versiones (CRID).
\end{itemize}

El \texttt{.h5} completo solo hace falta para \path{pipeline.etapa1} (inspección opcional, §\ref{sec:inspeccion}) o para trabajar sin
internet con \texttt{-{}-h5-local}. También se puede compartir la carpeta \path{data/raw/} (\textasciitilde 460\,MB) por fuera de git; cada
\path{scene.json} conserva su procedencia.

\section{Flujo}

\begin{lstlisting}
config.yaml --> pipeline.catalog --> data/catalog/*.json
                                          |
                                          v
                  .h5 (ASF, ~7 GB) --> pipeline.download --> data/raw/<scene_id>/  --> P2
                                          |                    (5 .tif + scene.json)
                                          +--> comparación del par (COMPARABLE / NO COMPARABLE)

.h5 local --> pipeline.etapa1 --> data/explore/<granuleId>/   (opcional, inspección)
\end{lstlisting}

Diagrama detallado: artefacto «Flujo de datos P1» (privado, compartido desde la página).

\subsection{Catálogo: \texttt{python -m pipeline.catalog}}

No requiere \texttt{.netrc}. Tarda menos de un minuto.

\begin{enumerate}
  \item Lee \path{roi.wkt} y \path{search.product} de \path{config.yaml}. \path{search.start}/\texttt{end} quedan en \texttt{null} a propósito para ver todas las pilas.
  \item Busca escenas NISAR del producto que intersectan la ROI.
  \item Repite la búsqueda filtrando por cada madurez (BETA, PROVISIONAL, VALIDATED), porque \path{asf_search} no la devuelve como
        propiedad. Una escena que no aparece en ningún filtro queda con madurez \texttt{None} y \textbf{no será comparable}.
  \item Convierte cada resultado en \texttt{SceneMeta} sin rellenar campos que falten.
  \item Agrupa en pilas por (producto, trayectoria, marco, dirección, polarizaciones, CRID, madurez).
  \item Imprime las pilas (\texttt{CANDIDATA} si tienen 2 o más fechas) y guarda
        \path{data/catalog/catalog_<fuente>_<producto>_<UTC>.json} con la consulta, la versión del código y todas las escenas.
\end{enumerate}

\subsection{Recorte: \texttt{python -m pipeline.download}}

\begin{lstlisting}
python -m pipeline.download <scene_id> [<scene_id> ...] [opciones]
\end{lstlisting}

\begin{longtable}{p{0.38\linewidth}p{0.55\linewidth}}
\toprule
\textbf{Opción} & \textbf{Uso} \\
\midrule
\endhead
\path{--capas HHHH HVHV mask ...} & Capas a recortar. Por defecto: \path{HHHH HVHV mask numberOfLooks rtcGammaToSigmaFactor} \\
\path{--h5-local "ruta.h5"} & Usar un \texttt{.h5} ya descargado (solo una escena) \\
\path{--catalogo ruta.json} & Usar un catálogo concreto; por defecto, el más reciente de \path{data/catalog/} \\
\path{--config ruta.yaml} & Otra configuración \\
\bottomrule
\end{longtable}

Por cada escena:

\begin{enumerate}
  \item \textbf{Busca la escena} en el catálogo. Si no está, se detiene: cada recorte queda ligado a una consulta guardada.
  \item \textbf{Abre el \texttt{.h5}.} En remoto: carga \texttt{.env}, inicia sesión en Earthdata
        (\path{EARTHDATA_USERNAME}/\path{EARTHDATA_PASSWORD} si están, y si no \texttt{.netrc}) y lo abre con \texttt{h5py} sobre HTTPS
        (\texttt{fsspec}, bloques de 8\,MiB con caché). En local: lo abre del disco.
  \item \textbf{Verifica la identidad.} \texttt{granuleId}, \texttt{compositeReleaseId}, \texttt{trackNumber}, \texttt{frameNumber} y
        \texttt{orbitPassDirection} del archivo deben coincidir con el catálogo. Si no coinciden, se detiene.
  \item \textbf{Lee la malla} de \path{science/LSAR/GCOV/grids/frequencyA}: EPSG, coordenadas, espaciado, polarizaciones, capas de
        potencia (términos diagonales de la covarianza) y su magnitud (\texttt{gamma0}, leída del archivo).
  \item \textbf{Calcula la ventana}: densifica la ROI, la transforma al CRS de la malla y toma el rectángulo de filas y columnas que la contiene.
  \item \textbf{Lee cada capa} solo en esa ventana, por franjas alineadas a los chunks del HDF5. En remoto, solo se transfieren esos
        bloques (\textasciitilde 230\,MB por escena en lugar de \textasciitilde 7\,GB).
  \item \textbf{Escribe un GeoTIFF por capa}: CRS de la escena, la transformación de la ventana, LZW, teselas de 512. Primero escribe
        \texttt{.tif.part} y luego lo renombra, así un corte no deja archivos a medias.
  \item \textbf{Mide la cobertura} con la capa \texttt{mask} dentro del \textbf{polígono} de la ROI (no del rectángulo).
  \item \textbf{Escribe o actualiza \texttt{scene.json}.} Si ya existía con otra malla o ventana (por ejemplo, porque cambió la ROI), se detiene.
\end{enumerate}

Con dos escenas o más, al final compara cada una contra la primera:

\begin{itemize}
  \item \path{compat.check_pair} (catálogo): producto, trayectoria, marco, dirección, polarizaciones, CRID y madurez iguales y conocidos.
  \item \path{subset.compare_grids} (\texttt{.h5}): frecuencia, EPSG, espaciado, forma y transformación del recorte idénticos.
\end{itemize}

Imprime \texttt{COMPARABLE} (código de salida 0) o \texttt{NO COMPARABLE} con los motivos (código 2).

Volver a ejecutar sobre la misma escena reescribe las capas y actualiza \path{scene.json} sin duplicar nada.

\subsection{Inspección: \texttt{python -m pipeline.etapa1}}
\label{sec:inspeccion}

\begin{lstlisting}
python -m pipeline.etapa1 "H5 FILES\archivo.h5" [--sin-vista] [--sin-figuras]
\end{lstlisting}

Solo con archivos locales. Escribe en \path{data/explore/<granuleId>/}: \path{metadatos.json}, \path{LEEME_validacion.txt}
(comprobaciones a mano con sus valores), \path{geotiff/} (\gz, \sz, dB con multilook 2\,\por\,2, máscara, factor, looks),
\path{figuras/} y \path{control/} (píxeles y bloques elegidos con semilla fija). \texttt{-{}-sin-vista} omite la vista general, que
obliga a leer la escena completa.

\section{Contrato C1 (lo que recibe P2)}

\begin{lstlisting}
data/raw/<scene_id>/
|-- HHHH.tif                  gamma0 HH, potencia lineal, float32, nodata NaN
|-- HVHV.tif                  gamma0 HV, potencia lineal, float32, nodata NaN
|-- mask.tif                  uint8 sin nodata: 0 = no válido, 255 = fuera de la adquisición, otro = subfranja válida
|-- numberOfLooks.tif         float32, nodata NaN
|-- rtcGammaToSigmaFactor.tif float32, nodata NaN   (sigma0 = gamma0 * factor)
`-- scene.json
\end{lstlisting}

Todas las capas comparten CRS, transformación y forma. Los valores son los del producto: \textbf{sin dB, sin máscara aplicada y sin
multilook}, en el rectángulo que contiene la ROI (el polígono se aplica en el procesamiento).

\subsection*{Campos de \texttt{scene.json}}

{\small
\begin{longtable}{p{0.20\linewidth}p{0.73\linewidth}}
\toprule
\textbf{Campo} & \textbf{Contenido} \\
\midrule
\endhead
\texttt{contrato} & \texttt{"C1"} \\
\texttt{scene} & \texttt{SceneMeta} del catálogo: \path{scene_id}, \texttt{product}, \texttt{track}, \texttt{frame}, \texttt{direction},
  \texttt{polarizations}, \texttt{crid}, \texttt{maturity}, \path{start_time}, \path{frame_coverage}, \path{orbit_type},
  \path{product_version}, \texttt{url}, \path{size_bytes}, \texttt{extra} \\
\path{identificacion_h5} & Lo que dice el propio archivo: \texttt{granuleId}, \texttt{compositeReleaseId}, \texttt{trackNumber},
  \texttt{frameNumber}, \texttt{orbitPassDirection}, \texttt{absoluteOrbitNumber}, \texttt{zeroDopplerStartTime}/\texttt{EndTime},
  \texttt{productVersion}, \texttt{isFullFrame}, \texttt{listOfFrequencies} \\
\texttt{roi} & \texttt{wkt} y \texttt{crs} usados para el recorte \\
\texttt{malla} & \texttt{frecuencia}, \texttt{epsg}, \path{dx_m}, \path{dy_m}, \path{forma_malla_completa}, \path{magnitud_potencia},
  \texttt{polarizaciones}, \path{capas_potencia}, \texttt{ventana} (\texttt{fila0}, \texttt{fila1}, \texttt{col0}, \texttt{col1}),
  \path{forma_recorte}, \path{transform_recorte} (6 coeficientes afines), \path{roi_se_sale_de_la_escena} \\
\texttt{capas} & Por capa: \texttt{archivo}, \texttt{descripcion}, \texttt{dtype}, \texttt{nodata}, \texttt{unidades} (del atributo
  del \texttt{.h5}), \path{creado_utc} \\
\path{cobertura_roi} & \path{pixeles_en_poligono}, \path{pct_valido}, \path{pct_no_valido}, \path{pct_fuera_adquisicion}, \texttt{nota} \\
\texttt{procedencia} & \texttt{fuente}, \texttt{metodo}, y \texttt{url} + \path{block_size_bytes} (remoto) o \path{ruta_local} +
  \path{url_catalogo} (local) \\
\path{code_version} & Commit corto de git (\texttt{-dirty} si había cambios sin commit) \\
\path{actualizado_utc} & Última escritura \\
\texttt{advertencias} & Recordatorios para quien consuma el recorte \\
\bottomrule
\end{longtable}
}

\textbf{Cambio respecto al plan original:} C1 pedía «el comando GDAL usado». Se reemplazó por \texttt{procedencia}, porque el recorte
se hace con \texttt{h5py}: el GDAL del entorno (3.12) no trae los drivers HDF5 ni NISAR.

\section{Resultados}

Catálogo del 2026-11-14 (\path{data/catalog/catalog_asf_GCOV_20261114T231456Z.json}): 39 escenas GCOV sobre la ROI.

\textbf{Par de trabajo:} GCOV T62/F10 ASCENDING, HH+HV, P05023, PROVISIONAL.

\begin{longtable}{p{0.22\linewidth}p{0.35\linewidth}p{0.35\linewidth}}
\toprule
 & \textbf{t1} & \textbf{t2} \\
\midrule
\endhead
Fecha & 2026-06-18 & 2026-10-16 \\
\path{scene_id} & \path{NISAR_L2_PR_GCOV_023_062_A_010_4005_DHDH_A_20260618T114010_20260618T114044_P05023_N_F_J_001}
                & \path{NISAR_L2_PR_GCOV_033_062_A_010_4005_DHDH_A_20261016T114006_20261016T114041_P05023_N_F_J_001} \\
Píxeles válidos en el polígono ROI & 99.57\,\% & 99.61\,\% \\
Comparabilidad & \texttt{COMPARABLE} & \\
\bottomrule
\end{longtable}

\begin{itemize}
  \item Malla: EPSG:32615 (UTM 15N), 10\,\por\,10\,m; malla completa de 33480\,\por\,34128\,px; recorte de 2687\,\por\,5504\,px,
        idéntico en ambas fechas.
  \item Tamaño en disco: \textasciitilde 230\,MB por escena (5 capas). Tiempo: \textasciitilde 3 minutos por par.
\end{itemize}

\textbf{Otras escenas recortadas:}

\begin{longtable}{p{0.30\linewidth}p{0.38\linewidth}p{0.24\linewidth}}
\toprule
\textbf{Escenas} & \textbf{Cobertura de la ROI} & \textbf{Uso} \\
\midrule
\endhead
T55/F80 DESC, P05023, 2026-06-18 y 2026-10-16 & \textasciitilde 50\,\% (32.6\,\% fuera de la adquisición, 17\,\% no válido; datos solo
  al este de lon \aprox{} $-93.14$) & Respaldo. Es COMPARABLE, pero cubre la mitad \\
T62/F10 ASC, P05036, 2026-10-28 & 99.6\,\% & Sin pareja: única fecha P05036 de su pila \\
\bottomrule
\end{longtable}

\textbf{Verificación con datos reales.} El recorte de la escena del 2026-10-28 desde su \texttt{.h5} local se comparó contra la salida
independiente de \path{etapa1.py}: misma transformación, misma forma y valores idénticos en los 14\,699\,442 píxeles válidos.

\section{Decisiones tomadas}

\begin{longtable}{p{0.36\linewidth}p{0.57\linewidth}}
\toprule
\textbf{Decisión} & \textbf{Motivo} \\
\midrule
\endhead
Recortar con \texttt{h5py} en lugar de \texttt{gdalwarp} & El GDAL 3.12 de conda-forge no trae los drivers HDF5 ni NISAR.
  \texttt{h5py} ya estaba en el entorno: no se agregaron dependencias \\
Leer el \texttt{.h5} remoto por ventanas & Evita descargar \textasciitilde 7\,GB por escena (\texttt{AGENTS.md} §5) \\
Par T62 en lugar de T55 & T55 solo cubre \textasciitilde 50\,\% de la ROI \\
No aplicar dB, máscara ni multilook en el recorte & Separar adquisición (P1) de procesamiento (P2); el recorte se puede comparar
  contra el producto valor por valor \\
Verificar la identidad del archivo contra el catálogo & Un recorte con la identidad equivocada rompe la trazabilidad (\texttt{AGENTS.md} §2.3) \\
Medir la cobertura sobre el polígono & El rectángulo incluye zonas fuera de la ROI \\
No crear \path{sources/event.py} & Los recursos del reto para NISAR son ASF/Earthdata, los mismos que ya usa \texttt{ASFSource} \\
\path{search.start}/\texttt{end} en \texttt{null} & Que el catálogo siga mostrando todas las pilas (por ejemplo, fechas P05036 nuevas) \\
\bottomrule
\end{longtable}

\section{Pruebas}
\label{sec:pruebas}

\begin{lstlisting}
pytest
\end{lstlisting}

Todas usan HDF5 \textbf{sintéticos} y no necesitan red. Las de P1 son 28 casos (de 44 en total):

\begin{longtable}{p{0.20\linewidth}p{0.12\linewidth}p{0.60\linewidth}}
\toprule
\textbf{Archivo} & \textbf{Funciones de prueba} & \textbf{Qué cubre} \\
\midrule
\endhead
\path{test_sources.py} & 4 & Normalización de resultados de ASF a \texttt{SceneMeta} \\
\path{test_nisar_h5.py} & 17 & Malla, ventana de la ROI, lectura por franjas, potencia, dB, multilook, estadísticas \\
\path{test_subset.py} & 7 & El recorte copia valores exactos y la georreferencia; cobertura sobre el polígono; rechazo por identidad
  y por capa inexistente; acumulación de capas en \path{scene.json}; detección de mallas desplazadas; ruta local de \texttt{ASFSource} \\
\bottomrule
\end{longtable}

\section{Problemas frecuentes}

\begin{longtable}{p{0.40\linewidth}p{0.53\linewidth}}
\toprule
\textbf{Error} & \textbf{Solución} \\
\midrule
\endhead
\texttt{No se pudo iniciar sesión en Earthdata} & Revisa \texttt{\%USERPROFILE\%\textbackslash .netrc} y las credenciales; el mensaje
  incluye el motivo de Earthdata \\
\texttt{No están en catalog\_...json} & Corre \path{python -m pipeline.catalog} o revisa que el \path{scene_id} esté completo \\
\texttt{El archivo no coincide con la escena del catálogo} & El \texttt{.h5} no es la escena pedida (revisa el ID o \texttt{-{}-h5-local}) \\
\texttt{ya tiene otra malla/ventana} & Cambió la ROI: borra \path{data/raw/<scene_id>/} y vuelve a recortar \\
\texttt{Capas no disponibles en el archivo} & Esa escena no tiene la capa (por ejemplo, una escena solo HV); revisa \texttt{-{}-capas} \\
\texttt{ModuleNotFoundError: pipeline} & Ejecuta desde \texttt{C:\textbackslash terrapulse} con \path{python -m pipeline.<script>} \\
\texttt{UnicodeEncodeError} en la consola de Windows & Define \path{set PYTHONIOENCODING=utf-8} antes de ejecutar \\
\bottomrule
\end{longtable}

\section{Limitaciones y pendientes}

\begin{itemize}
  \item \textbf{\texttt{code\_version} de los recortes:} el par T62 se regeneró el 2026-11-15 con el árbol limpio y registra \texttt{19de237}.
        Los recortes de respaldo (T55) y la escena P05036 conservan versiones anteriores con \texttt{-dirty}; regenerarlos si se van a usar.
  \item \textbf{Mínimo de cobertura:} lo fija P3.
  \item \textbf{Línea base para el z-score (Fase 2):} recortar más fechas de la pila T62 P05023 con \path{pipeline.download <scene_id>}
        (\textasciitilde 1.5\,min por fecha), cuando P2 diga cuántas necesita.
  \item \textbf{Pila T156/F80:} su cobertura de la ROI no se ha verificado.
  \item \textbf{Antes del evento:} volver a correr \path{pipeline.catalog} por si aparecen fechas nuevas (por ejemplo, T62 P05036 para
        emparejar la escena del 2026-10-28).
  \item \textbf{Advertencias del producto} que P2 debe atender: sin corrección de ruido (\path{noiseCorrectionApplied = False}), sin
        máscara de sombra/layover (\path{shadowMaskingApplied = False}), potencias con mantisa de 16 bits (\textasciitilde 5 cifras
        significativas). Productos PROVISIONAL: calibrados y parcialmente validados, no completamente validados.
  \item \path{etapa1.py} solo trabaja con archivos locales.
\end{itemize}

\end{document}
